\section{Limitations}
\label{sec:limitations}

\paragraph{Scope and evaluation limits.}
The prolonged residual $R$ is evaluated against a procedural no-cliff reference and does not capture human perceptual nuances directly.
Formal listening evaluations (MOS, MUSHRA) require controlled subjective listening panels and are deferred to dedicated perceptual studies.
Speech-oriented objective metrics (PESQ, STOI) are excluded because the evaluation cycles consist of non-speech synthetic wavetable tiles; applying speech intelligibility models to elementary harmonic and multi-family geometric tiles represents an ecological domain mismatch.

\begin{itemize}
  \item \textbf{Reference geometry.} Residual scores quantify agreement against a procedural ideal sibling rather than subjective acoustic preferences.
  \item \textbf{Transfer domain.} The primary overnight search operates on Python \texttt{make\_batch} multi-family draws; the Rust \texttt{sound\_bench} suite serves as an out-of-distribution transfer benchmark.
  \item \textbf{Search bounds.} Per-trial architecture depth and width are bounded to preserve overnight search throughput ($5{,}000$ iterations on a single GPU). Diffusion models and full-scale Demucs variants were screened out.
  \item \textbf{Freeze vs.\ re-run protocols.} The $5$k-gate freeze captures organic branch co-evolution, whereas ablation re-runs test isolated mechanics over constrained budgets under the identical pseudo-random seed.
  \item \textbf{Classical baselines.} VA and bandlimited BLEP implementations follow canonical formulations from open-access literature~\cite{stilson1996,nam2009polyblep,esqueda2016blamp}.
\end{itemize}


